\documentclass[../../main.tex]{subfiles}

\begin{document}

\section{Arzel\`a-Ascoli 定理}\label{section:ch10-arzela-ascoli}

在分析学的许多重要问题中, 给定连续实值函数序列, 知道是否有一致收敛的子序列是有用的. 本节我们的主要结果是 Arzel\`a-Ascoli 定理, 它提供了紧度量空间 $X$ 上的一致有界的连续实值函数序列有一致收敛的子序列的准则. 在证明该定理后, 我们将它与找出度量空间的子集为紧的判别准则的一般问题相联系.

对于一个度量空间 $X$, 我们将 $X$ 上的连续实值函数的线性空间记为 $C(X)$. 若 $X$ 是紧的, 根据前一章的定理 22, $X$ 上的每个连续函数取到最大值. 对于 $C(X)$ 中的函数 $f$, 定义
\begin{align*}
\|f\|_{\max}=\max_{x\in X}|f(x)|.
\end{align*}
如同我们首先在第 7 章考虑过的特殊情形 $X=[a,b]$ 那样, 这定义了一个范数. 这个最大值范数诱导了度量
\begin{align*}
\rho_{\max}(g,h)=\|g-h\|_{\max},\quad\text{对所有 }g,h\in C(X).
\end{align*}
我们称该度量为一致度量, 因为 $C(X)$ 中的序列关于该度量收敛当且仅当它在 $X$ 上一致收敛. 关于该度量的 Cauchy 序列称为一致 Cauchy 序列. $C(X)$ 的完备性的证明与 $X=[a,b]$ 的情形的证明并无不同(见前一章的命题 10 的证明).

\begin{proposition}\label{proposition:ch10-CX完备}
若 $X$ 是紧度量空间, 则 $C(X)$ 是完备的.
\end{proposition}

\begin{definition}[等度连续]\label{definition:ch10-等度连续}
度量空间 $X$ 上的实值函数族 $\mathcal{F}$ 称为在点 $x\in X$ 等度连续, 若对每个 $\varepsilon>0$, 存在 $\delta>0$ 使得对每个 $f\in\mathcal{F}$ 和 $x'\in X$,
\begin{align*}
\text{若 }\rho(x',x)<\delta,\quad\text{则 }|f(x')-f(x)|<\varepsilon.
\end{align*}
族 $\mathcal{F}$ 称为在 $X$ 上等度连续, 若它在 $X$ 的每个点等度连续.
\end{definition}

当然, 等度连续族中的每个函数是连续的, 而任何有限多个连续函数族是等度连续的. 一般说来, 无穷多个连续函数的族不是等度连续的. 例如, 对每个自然数 $n$, 定义 $f_{n}(x)=x^{n}, 0\leqslant x\leqslant 1$, 则 $\{f_{n}\}$ 是 $[0,1]$ 上的连续函数的可数族, 它在 $x=0$ 不是等度连续的, 而在 $[0,1]$ 的其余点是等度连续的.

\begin{example}\label{example:ch10-等度连续例}
对 $M\geqslant 0$, 令 $\mathcal{F}$ 为在闭有界区间 $[a,b]$ 上连续, 在开区间 $(a,b)$ 上可微且在 $(a,b)$ 上 $|f'|\leqslant M$ 的函数族. 我们从中值定理推出对所有 $u,v\in[a,b]$,
\begin{align*}
|f(u)-f(v)|\leqslant M\cdot|u-v|.
\end{align*}
因此 $\mathcal{F}$ 是等度连续的, 由于考虑到在 $X$ 的每一点等度连续的准则, $\delta=\varepsilon/M$ 对应于 $\varepsilon>0$ 的挑战.
\end{example}

集合 $X$ 上的实值函数序列 $\{f_{n}\}$ 称为逐点有界, 若对每个 $x\in X$, 序列 $\{f_{n}(x)\}$ 是有界的; 而称为在 $X$ 上一致有界, 若存在某个 $M\geqslant 0$ 使得对所有 $n$, 在 $X$ 上 $|f_{n}|\leqslant M$.

\begin{lemma}[Arzel\`a-Ascoli 引理]\label{lemma:ch10-arzela-ascoli引理}
令 $X$ 为可分度量空间, 而 $\{f_{n}\}$ 是 $C(X)$ 中的逐点有界的等度连续函数序列. 则 $\{f_{n}\}$ 的一个子序列在整个 $X$ 上逐点收敛到 $X$ 上的实值函数 $f$.
\end{lemma}

\begin{proof}
令 $\{x_{j}\}_{j=1}^{\infty}$ 为 $X$ 的稠密子集 $D$ 的列举. 定义为 $n\mapsto f_{n}(x_{1})$ 的实数序列是有界的. 因此, 根据 Bolzano-Weierstrass 定理, 该序列有收敛的子序列, 即存在严格递增的整数序列 $\{s(1,n)\}$ 以及数 $a_{1}$ 使得
\begin{align*}
\lim_{n\to\infty}f_{s(1,n)}(x_{1})=a_{1}.
\end{align*}
用同样的方法, 定义为 $n\mapsto f_{s(1,n)}(x_{2})$ 的序列是有界的, 且因此存在一个 $\{s(1,n)\}$ 的子序列 $\{s(2,n)\}$ 和数 $a_{2}$ 使得 $\lim_{n\to\infty}f_{s(2,n)}(x_{2})=a_{2}$. 我们归纳地继续这个选取过程以得到一个严格递增的实数序列 $\{\{s(j,n)\}\}_{j=1}^{\infty}$ 和一个数序列 $\{a_{j}\}$, 使得对每个指标 $j$,
\begin{align*}
\{s(j+1,n)\}\text{ 是 }\{s(j,n)\}\text{ 的子序列且 }\lim_{n\to\infty}f_{s(j,n)}(x_{j})=a_{j}.
\end{align*}
对于每个指标 $j$, 定义 $f(x_{j})=a_{j}$. 考虑通过对每个指标 $k$ 设 $n_{k}=s(k,k)$ 得到的"对角"序列 $\{f_{n_{k}}\}$. 对每个 $j$, $\{n_{k}\}_{k=j}^{\infty}$ 是上面所选的自然数的 $j$ 重子序列, 于是,
\begin{align*}
\lim_{k\to\infty}f_{n_{k}}(x_{j})=a_{j}=f(x_{j}).
\end{align*}
因此 $\{f_{n_{k}}\}$ 在 $D$ 上逐点收敛到 $f$.

为了记号上的方便, 假设整个序列 $\{f_{n}\}$ 在 $D$ 上逐点收敛到 $f$. 令 $x_{0}$ 为 $X$ 中的任意一点. 我们宣称 $\{f_{n}(x_{0})\}$ 是 Cauchy 的. 事实上, 令 $\varepsilon>0$. 根据 $\{f_{n}\}$ 在 $x_{0}$ 点的等度连续性, 我们可以选取 $\delta>0$, 使得 $|f_{n}(x)-f_{n}(x_{0})|<\varepsilon/3$, 对所有指标 $n$ 和所有满足 $\rho(x,x_{0})<\delta$ 的 $x\in X$ 成立. 由于 $D$ 是稠密的, 存在点 $x\in D$ 使得 $\rho(x,x_{0})<\delta$. 此外, 由于 $\{f_{n}(x)\}$ 收敛, 它必须是一个 Cauchy 序列, 因而我们可以选取 $N$ 充分大, 使得对所有 $m,n\geqslant N$,
\begin{align*}
|f_{n}(x)-f_{m}(x)|<\varepsilon/3.
\end{align*}
则对所有 $m,n\geqslant N$,
\begin{align*}
|f_{n}(x_{0})-f_{m}(x_{0})|&\leqslant|f_{n}(x_{0})-f_{n}(x)|+|f_{n}(x)-f_{m}(x)|+|f_{m}(x_{0})-f_{m}(x)|\\
&<\varepsilon/3+\varepsilon/3+\varepsilon/3=\varepsilon.
\end{align*}
因此 $\{f_{n}(x_{0})\}$ 是一个实数的 Cauchy 序列. 由于 $\mathbf{R}$ 是完备的, $\{f_{n}(x_{0})\}$ 收敛. 将该极限记为 $f(x_{0})$. 序列 $\{f_{n}\}$ 在整个 $X$ 上逐点收敛到 $f:X\to\mathbf{R}$.
\end{proof}

我们证明了紧度量空间上的连续实值函数是一致连续的. 完全同样的证明表明, 若 $X$ 是一个紧度量空间, 而 $\mathcal{F}$ 是 $X$ 上的等度连续函数族, 则 $\mathcal{F}$ 在以下意义下是一致等度连续的: 对每个 $\varepsilon>0$, 存在 $\delta>0$ 使得对 $u,v\in X$ 和任何 $f\in\mathcal{F}$,
\begin{align*}
\text{若 }\rho(u,v)<\delta,\quad\text{则 }|f(u)-f(v)|<\varepsilon.
\end{align*}

\begin{theorem}[Arzel\`a-Ascoli 定理]\label{theorem:ch10-arzela-ascoli}
令 $X$ 为紧度量空间, 而 $\{f_{n}\}$ 是 $X$ 上的一致有界的等度连续的实值函数序列. 则 $\{f_{n}\}$ 有一个在 $X$ 上一致收敛于 $X$ 上的连续函数 $f$ 的子序列.
\end{theorem}

\begin{proof}
由于 $X$ 是紧度量空间, 根据前一章的命题 24, 它是可分的. Arzel\`a-Ascoli 引理告诉我们 $\{f_{n}\}$ 的子序列在整个 $X$ 上逐点收敛于实值函数 $f$. 为了记号上的方便, 假设整个序列 $\{f_{n}\}$ 在 $X$ 上逐点收敛. 因此, 特别地, 对 $X$ 中的每个 $x$, $\{f_{n}(x)\}$ 是实数的 Cauchy 序列. 我们用这个结论和等度连续性证明 $\{f_{n}\}$ 是 $C(X)$ 中的 Cauchy 序列.

令 $\varepsilon>0$. 根据 $\{f_{n}\}$ 在 $X$ 上的一致等度连续性, 存在 $\delta>0$ 使得对所有 $n$, 以及所有使得 $\rho(u,v)<\delta$ 的 $u,v\in X$,
\begin{align}
|f_{n}(u)-f_{n}(v)|<\varepsilon/3.\label{align:ch10-arzela-1}
\end{align}
由于 $X$ 是紧度量空间, 根据前一章的定理 16, 它是全有界的. 因此存在 $X$ 中的有限个点 $x_{1},\cdots,x_{k}$ 使得 $X$ 被 $\{B(x_{i},\delta)\}_{i=1}^{k}$ 覆盖. 对 $1\leqslant i\leqslant k$, $\{f_{n}(x_{i})\}$ 是 Cauchy 的, 因此存在 $N$ 使得对 $1\leqslant i\leqslant k$ 与所有 $n,m\geqslant N$,
\begin{align}
|f_{n}(x_{i})-f_{m}(x_{i})|<\varepsilon/3.\label{align:ch10-arzela-2}
\end{align}
现在对 $X$ 中的任何 $x$, 存在 $i$, $1\leqslant i\leqslant k$, 使得 $\rho(x,x_{i})<\delta$, 因此对 $n,m\geqslant N$,
\begin{align*}
|f_{n}(x)-f_{m}(x)|&\leqslant|f_{n}(x)-f_{n}(x_{i})|+|f_{n}(x_{i})-f_{m}(x_{i})|+|f_{m}(x_{i})-f_{m}(x)|\\
&<\varepsilon/3+\varepsilon/3+\varepsilon/3=\varepsilon.
\end{align*}
因此 $\{f_{n}\}$ 是一致 Cauchy 的. 于是, 由于 $C(X)$ 是完备的, $\{f_{n}\}$ 在 $X$ 上一致收敛到一个连续函数.
\end{proof}

我们证明了一个度量空间是紧的当且仅当它是序列紧的. 此外, 对于欧氏空间 $\mathbf{R}^{n}$ 的子空间 $K$, Heine-Borel 定理告诉我们 $K$ 是紧的当且仅当 $K$ 是 $\mathbf{R}^{n}$ 的闭有界子集. 在一般的度量空间, 闭且有界是紧的必要条件, 但不是充分条件. 例如, $C[0,1]$ 的闭单位球 $\{f\in C[0,1]\mid\|f\|_{\max}\leqslant 1\}$ 是 $C[0,1]$ 的闭有界子集但不是序列紧的. 的确, $[0,1]$ 上对所有 $n$ 定义为 $f_{n}(x)=x^{n}$ 的序列 $\{f_{n}\}$ 没有一致收敛到 $[0,1]$ 上连续函数的子序列. Arzel\`a-Ascoli 定理可作为确定 $C(X)$ 的闭有界子集是紧的准则来阐述.

\begin{theorem}\label{theorem:ch10-CX紧集准则}
令 $X$ 为紧度量空间而 $\mathcal{F}$ 是 $C(X)$ 的子集, 则 $\mathcal{F}$ 是 $C(X)$ 的紧子集当且仅当 $\mathcal{F}$ 是闭的、一致有界且等度连续.
\end{theorem}

\begin{proof}
首先假定 $\mathcal{F}$ 是闭的、一致有界且等度连续. 令 $\{f_{n}\}$ 为 $\mathcal{F}$ 中的序列. 根据 Arzel\`a-Ascoli 定理, $\{f_{n}\}$ 的一个子序列收敛到函数 $f\in C(X)$. 由于 $\mathcal{F}$ 是闭的, 所以 $f$ 属于 $\mathcal{F}$. 因此 $\mathcal{F}$ 是序列紧度量空间, 因此是紧的.

现在假设 $\mathcal{F}$ 是紧的. 我们把证明 $\mathcal{F}$ 是 $C(X)$ 的有界闭子集留作练习. 我们用反证法证明 $\mathcal{F}$ 是等度连续的. 假定 $\mathcal{F}$ 在 $X$ 中的点 $x$ 处不是等度连续的. 则存在 $\varepsilon_{0}>0$ 使得对每个自然数 $n$, 存在 $\mathcal{F}$ 中的函数 $f_{n}$ 和 $X$ 中的点 $x_{n}$ 使得
\begin{align}
|f_{n}(x_{n})-f_{n}(x)|\geqslant\varepsilon_{0}\quad\text{而 }\rho(x_{n},x)<1/n.\label{align:ch10-紧集-3}
\end{align}
由于 $\mathcal{F}$ 是紧度量空间, 它是序列紧的. 因此存在子序列 $\{f_{n_{k}}\}$ 在 $X$ 上一致收敛到连续函数 $f$. 选取一个指标 $K$ 使得对于 $k\geqslant K$, $\rho_{\max}(f,f_{n_{k}})<\varepsilon_{0}/3$. 我们从\eqref{align:ch10-紧集-3}推出对 $k\geqslant K$,
\begin{align}
|f(x_{n_{k}})-f(x)|>\varepsilon_{0}/3\quad\text{而 }\rho(x_{n_{k}},x)<1/n_{k}.\label{align:ch10-紧集-4}
\end{align}
这与 $f$ 在点 $x$ 连续矛盾. 因此 $\mathcal{F}$ 是等度连续的.
\end{proof}

\begin{remark}\label{remark:ch10-arzela-注}
Arzel\`a-Ascoli 引理的证明非常类似于第 7 章的 Helly 定理的证明. 构成这两个证明的基础的共同技巧是所谓的 Cantor 对角线方法(见习题 14).
\end{remark}

\begin{remark}\label{remark:ch10-lp紧}
第 13 章的 Riesz 定理告诉我们, 赋范线性空间的闭单位球是紧的当且仅当线性空间是有限维的. 因此, 给定一个特定的无穷维赋范线性空间, 刻画那些紧的闭有界子集是有趣的. Arzel\`a-Ascoli 定理给出的关于 $C(X)$ 的子空间的紧性判别准则在 $\ell^{p}$ 有对应的结果. 不难证明, 对于 $1\leqslant p<\infty$, $\ell^{p}$ 的闭有界子集是紧的当且仅当它在对每个 $\varepsilon>0$, 存在指标 $N$ 使得
\begin{align*}
\text{对所有 }x=\{x_{n}\}\in S,\quad\sum_{k=N}^{\infty}|x_{k}|^{p}<\varepsilon
\end{align*}
的意义下是等度可和的.
\end{remark}

\begin{exercise}\label{exercise:ch10-01}
令 $E$ 为度量空间 $Y$ 的紧子空间. 证明 $E$ 是 $Y$ 的闭有界子集.
\end{exercise}

\begin{exercise}\label{exercise:ch10-02}
证明紧度量空间上等度连续的实值函数序列是逐点有界的当且仅当它是一致有界的.
\end{exercise}

\begin{exercise}\label{exercise:ch10-03}
证明紧度量空间上连续函数的等度连续族是一致等度连续的.
\end{exercise}

\begin{exercise}\label{exercise:ch10-04}
令 $X$ 为度量空间, $\{f_{n}\}$ 为 $C(X)$ 中在 $X$ 上一致收敛到 $f\in C(X)$ 的序列. 证明 $\{f_{n}\}$ 是等度连续的.
\end{exercise}

\begin{exercise}\label{exercise:ch10-05}
$[0,1]$ 上的实值函数 $f$ 称为 $\alpha$ 阶 H\"older 连续, 若存在常数 $C$ 使得
\begin{align*}
\text{对所有 }x,y\in[0,1],\quad|f(x)-f(y)|\leqslant C|x-y|^{\alpha}.
\end{align*}
定义 H\"older 范数
\begin{align*}
\|f\|_{\alpha}=\max\{|f(x)|+|f(x)-f(y)|/|x-y|^{\alpha}\mid x,y\in[0,1],x\neq y\}.
\end{align*}
证明对于 $0<\alpha\leqslant 1$, 满足 $\|f\|_{\alpha}\leqslant 1$ 的函数集(作为 $C[0,1]$ 的子集)有紧的闭包.
\end{exercise}

\begin{exercise}\label{exercise:ch10-06}
令 $X$ 为紧度量空间而 $\mathcal{F}$ 是 $C(X)$ 的子集. 证明 $\mathcal{F}$ 是等度连续的当且仅当它在 $C(X)$ 的闭包 $\overline{\mathcal{F}}$ 等度连续. $C(X)$ 的子集有紧闭包当且仅当它等度连续且一致有界.
\end{exercise}

\begin{exercise}\label{exercise:ch10-07}
对于闭有界区间 $[a,b]$, 令 $\{f_{n}\}$ 为 $C[a,b]$ 中的序列. 若 $\{f_{n}\}$ 是等度连续的, $\{f_{n}\}$ 必须有一个一致收敛的子序列吗? 若 $\{f_{n}\}$ 一致有界, $\{f_{n}\}$ 必须有一个一致收敛的子序列吗?
\end{exercise}

\begin{exercise}\label{exercise:ch10-08}
令 $X$ 为紧度量空间而 $Y$ 是一般的度量空间. 记 $C(X,Y)$ 为从 $X$ 到 $Y$ 的连续映射集合. 对 $C(X,Y)$ 中的序列叙述和证明 Arzel\`a-Ascoli 定理. 这里 $\{f_{n}\}$ 逐点有界的假设被对每个 $x\in X$, 集合 $\{f_{n}(x)\mid n\text{ 是自然数}\}$ 的闭包是 $Y$ 的紧子空间的假设代替.
\end{exercise}

\begin{exercise}\label{exercise:ch10-09}
令 $\{f_{n}\}$ 为 $\mathbf{R}$ 上的等度连续、一致有界的连续实值函数序列. 证明存在 $\{f_{n}\}$ 的子序列在 $\mathbf{R}$ 上逐点收敛于 $\mathbf{R}$ 上的某个连续函数, 且该收敛在 $\mathbf{R}$ 的每个有界子集是一致的.
\end{exercise}

\begin{exercise}\label{exercise:ch10-10}
对于 $1\leqslant p<\infty$, 证明 $\ell^{p}$ 的一个子空间是紧的当且仅当它是闭、有界与等度可和的.
\end{exercise}

\begin{exercise}\label{exercise:ch10-11}
对于非负实数序列 $\{c_{n}\}$, 令 $S$ 为 $\ell^{p}$ 的那些使得对所有 $n$, $|x_{n}|\leqslant c_{n}$ 的 $x=\{x_{n}\}\in\ell^{p}$ 组成的子集. 证明: 若 $\{c_{n}\}$ 属于 $\ell^{p}$, $S$ 是等度可和的.
\end{exercise}

\begin{exercise}\label{exercise:ch10-12}
对于 $1\leqslant p\leqslant\infty$, 证明 Banach 空间 $\ell^{p}$ 的闭单位球不是紧的.
\end{exercise}

\begin{exercise}\label{exercise:ch10-13}
对于 $1\leqslant p\leqslant\infty$, 证明 Banach 空间 $L^{p}[0,1]$ 的闭单位球不是紧的.
\end{exercise}

\begin{exercise}\label{exercise:ch10-14}
令 $S$ 为可数集而 $\{f_{n}\}$ 为 $S$ 上的逐点有界的实值函数序列. 证明 $\{f_{n}\}$ 有在 $S$ 上逐点收敛到实值函数的子序列.
\end{exercise}

\end{document}